\scanpage{43}
\mspage{2}
% L1–6
les formules
\[
\unsure{\zeta|\xi}+\zeta'=12\eta-4\zeta=\deleted{\unsure{6}}4\lambda\tag*{(7)}
\]
donc (5) s'écrit \deleted{$\eta+\eta'=2\lambda$} \inserted{(en divisant par 2)}
\[
\eta+\eta'=2\lambda\tag*{(8)}
\]
ou encore
\[
\eta'=5\eta-2\zeta\tag*{(\unsure{8|9} bis)}
\]
% L7–13
Notons que $\lambda$ provient de l'application rationnelle $P\longrightarrow\mathbb P^3$ définie par le système des cubiques qui passent par $S$, laquelle devient régulière sur $\widetilde P$ et définit une immersion $\widetilde P\hookrightarrow\mathbb P^3$ dont l'image est une surface cubique. (En particulier, $\lambda$ est ample\ldots) \deleted{On}

% slanted
\margin{$\unsure{K|\lambda}=-3\unsure{\lambda|\eta}$ ($K$ classe canonique) ce qui caractérise $\lambda$ de façon intrinsèque, à la surface $\widetilde P$ $=X$. Du $=$ coup, l'inclusion $X\hookrightarrow\mathbb P^3$ (var. de \unsure{Sév.} Brauer) est intrinsèque à $X$, ainsi que les $D_i,D'_i,D_{ij}$, qui sont les seules droites de $\mathbb P^3$ contenues dans $X$ -- a fortiori, les 27 élts $\zeta_i,\zeta'_i,\zeta_{ij}\in E$ sont intrinsèques à $X$.}

Calculons la forme intersection sur $\widetilde P$.

On a
\[
\eta^2=1,\qquad \zeta_i^2=-1\qquad \eta\cdot\zeta_i=0,\quad\zeta_i\zeta_j=0\quad\text{si }i\ne j\tag*{(10)}
\]
d'où facilement
\[
\eta'{}^2=1,\quad\zeta'_i{}^2=-1,\quad\eta'\cdot\zeta'_i=0,\quad\zeta'_i\zeta'_j=0\quad\text{si }i\ne j.\tag*{(11)}
\]
\unsure{Et|et|On a} aussi
\[
\lambda^2=3\qquad(=\lambda\eta=\lambda\eta')\tag*{(12)}
\]
\[
\lambda\zeta_i=\lambda\zeta'_i=\lambda\zeta_{ij}=1\tag*{(13)}
\]
% L19–23
La relation (12) signifie que l'immersion $\widetilde P\longrightarrow\mathbb P^3$ a comme image une surface cubique, la relation (13) que les images dans $\mathbb P^3$ des 27 \inserted{$=6+6+15$} courbes $D_i,D'_i,D_{ij}$ sont des \uline{droites}.

Notons encore
\[
\zeta_i\zeta'_j=\begin{cases}0&\text{si }i=j\\1&\text{si }i\ne j\end{cases}\tag*{(14)}
\]
\unsure{(}
% archive sheet number: 42
